% LaTeX Template for AI Problem Analysis Output (v2.0)
% AMC 12B 2023 Problem 12 Strategic Analysis
% Structure your analysis results using this template with colorful boxes
% IMPORTANT: Use LaTeX commands for all mathematical symbols, NOT unicode characters

\documentclass[12pt]{article}
\usepackage[utf8]{inputenc}
\usepackage{amsmath, amssymb, amsthm}
\usepackage{geometry}
\usepackage{xcolor}
\usepackage[most]{tcolorbox}
\usepackage{enumitem}
\usepackage{fancyhdr}

% Page setup
\geometry{margin=1in}
\setlength{\headheight}{14.5pt}
\pagestyle{fancy}
\fancyhf{}
\rhead{AMC12 Problem Analysis}
\lhead{\today}

% Color definitions
\definecolor{problemconfig}{RGB}{230, 242, 255}    % Light blue
\definecolor{strategic}{RGB}{255, 230, 242}       % Light pink
\definecolor{tactical}{RGB}{242, 255, 230}        % Light green
\definecolor{techniques}{RGB}{255, 242, 230}      % Light yellow
\definecolor{verification}{RGB}{255, 230, 230}    % Light red
\definecolor{solution}{RGB}{230, 255, 255}        % Light cyan
\definecolor{competition}{RGB}{242, 242, 255}     % Light purple
\definecolor{proofwriting}{RGB}{240, 230, 255}    % Light lavender
\definecolor{gold}{RGB}{218, 165, 32}

% Box styles
\newtcolorbox{problemconfigbox}{
    colback=problemconfig,
    colframe=blue,
    title=Problem Configuration Analysis,
    fonttitle=\bfseries,
    arc=3pt,
    boxrule=1pt,
    breakable
}

\newtcolorbox{strategicbox}{
    colback=strategic,
    colframe=purple,
    title=Strategic Framework Application,
    fonttitle=\bfseries,
    arc=3pt,
    boxrule=1pt,
    breakable
}

\newtcolorbox{tacticalbox}{
    colback=tactical,
    colframe=green,
    title=Tactical Methodology Selection,
    fonttitle=\bfseries,
    arc=3pt,
    boxrule=1pt,
    breakable
}

\newtcolorbox{techniquesbox}{
    colback=techniques,
    colframe=orange,
    title=Conversion Techniques Application,
    fonttitle=\bfseries,
    arc=3pt,
    boxrule=1pt,
    breakable
}

\newtcolorbox{verificationbox}{
    colback=verification,
    colframe=red,
    title=Verification Strategy Design,
    fonttitle=\bfseries,
    arc=3pt,
    boxrule=1pt,
    breakable
}

\newtcolorbox{solutionbox}{
    colback=solution,
    colframe=cyan,
    title=Solution Roadmap,
    fonttitle=\bfseries,
    arc=3pt,
    boxrule=1pt,
    breakable
}

\newtcolorbox{competitionbox}{
    colback=competition,
    colframe=purple,
    title=Competition Strategy Integration,
    fonttitle=\bfseries,
    arc=3pt,
    boxrule=1pt,
    breakable
}

\newtcolorbox{keyinsightbox}{
    colback=yellow!20,
    colframe=gold,
    title=Key Insight,
    fonttitle=\bfseries,
    arc=3pt,
    boxrule=2pt,
    leftrule=5pt,
    breakable
}

\newtcolorbox{warningbox}{
    colback=red!10,
    colframe=red!50,
    title=Warning: Common Pitfall,
    fonttitle=\bfseries,
    arc=3pt,
    boxrule=2pt,
    leftrule=5pt,
    breakable
}

\newtcolorbox{tipbox}{
    colback=green!10,
    colframe=green!50,
    title=Pro Tip,
    fonttitle=\bfseries,
    arc=3pt,
    boxrule=2pt,
    leftrule=5pt,
    breakable
}

\begin{document}

\title{AMC12 Strategic Problem Analysis\\2023 AMC 12B Problem 12}
\author{Competition Math Framework v2.1}
\date{\today}
\maketitle

\section*{Problem Statement}

For complex number $u = a+bi$ and $v = c+di$ (where $i=\sqrt{-1}$), define the binary operation
$$u \otimes v = ac + bdi$$
Suppose $z$ is a complex number such that $z\otimes z = z^{2}+40$. What is $|z|$?

$$\textbf{(A) }2\qquad\textbf{(B) }5\qquad\textbf{(C) }\sqrt{5}\qquad\textbf{(D) }\sqrt{10}\qquad\textbf{(E) }5\sqrt{2}$$

\begin{problemconfigbox}
\subsection*{A. Problem Classification}
\begin{itemize}[leftmargin=*]
    \item \textbf{Problem Type}: To Find (compute the magnitude of a complex number)
    \item \textbf{Competition Context}: AMC12 2023B Problem 12
    \item \textbf{Difficulty Band}: Mid-band (P12/25) - requires understanding of complex numbers and algebraic manipulation
    \item \textbf{Mathematical Domain}: Complex Numbers / Algebra
    \item \textbf{Estimated Time}: 3-4 minutes for mid-band problem
\end{itemize}

\subsection*{B. Problem Structure Identification}
\begin{itemize}[leftmargin=*]
    \item \textbf{Given Information}: 
    \begin{itemize}
        \item Custom binary operation: $u \otimes v = ac + bdi$ for $u = a+bi$, $v = c+di$
        \item Condition: $z \otimes z = z^2 + 40$ for some complex number $z$
        \item Multiple choice answers available
    \end{itemize}
    \item \textbf{Constraints}: 
    \begin{itemize}
        \item $z$ is a complex number (can be written as $z = a+bi$ for real $a,b$)
        \item The custom operation $\otimes$ only takes the real part from first component ($ac$) and imaginary from second component ($bdi$)
    \end{itemize}
    \item \textbf{Objective}: Find $|z| = \sqrt{a^2+b^2}$
    \item \textbf{Hidden Assumptions}: 
    \begin{itemize}
        \item The operation $\otimes$ is not the standard complex multiplication
        \item We need to equate both real and imaginary parts separately
    \end{itemize}
    \item \textbf{Answer Choice Leverage}: All answers are positive real numbers, ranging from $2$ to $5\sqrt{2} \approx 7.07$. The presence of $\sqrt{10}$ and $5\sqrt{2} = \sqrt{50}$ suggests the magnitude might involve square roots of multiples of 10.
\end{itemize}

\subsection*{C. Strategic Orientation}
\begin{itemize}[leftmargin=*]
    \item \textbf{Hypothesis}: Let $z = a + bi$ where $a, b \in \mathbb{R}$
    \item \textbf{Conclusion}: Determine $|z| = \sqrt{a^2 + b^2}$
    \item \textbf{Notation}: 
    \begin{itemize}
        \item $z = a + bi$ (standard rectangular form)
        \item $z^2 = (a+bi)^2 = a^2 - b^2 + 2abi$
        \item $z \otimes z = a^2 + b^2i$ (applying the custom operation)
    \end{itemize}
    \item \textbf{Intuitive Approach}: 
    \begin{enumerate}
        \item Express $z$ in rectangular form
        \item Compute $z \otimes z$ using the given operation
        \item Compute $z^2$ using standard complex multiplication
        \item Set $z \otimes z = z^2 + 40$ and equate real/imaginary parts
        \item Solve for $a$ and $b$, then compute $|z|$
    \end{enumerate}
    \item \textbf{Cross-Domain Potential}: This is primarily an algebraic problem with complex numbers. Could also approach by testing answer choices (working backwards).
\end{itemize}
\end{problemconfigbox}

\begin{strategicbox}
\subsection*{A. Psychological Strategies}
\begin{itemize}[leftmargin=*]
    \item \textbf{Mental Toughness}: The custom operation may seem unusual, but remain calm and methodically apply it. Don't be intimidated by the non-standard notation.
    \item \textbf{Creativity Cultivation}: Recognize that this problem tests understanding of complex numbers beyond standard operations. The key is careful algebraic manipulation.
    \item \textbf{Iterative Refinement}: If the algebra becomes messy, step back and verify each computation. Check that you're applying the custom operation correctly (not standard multiplication).
\end{itemize}

\subsection*{B. Investigation Strategies}
\begin{itemize}[leftmargin=*]
    \item \textbf{Startup Orientation}: Read the custom operation definition carefully. Notice it's NOT standard complex multiplication - it mixes parts differently.
    \item \textbf{Get Your Hands Dirty}: Write out $z = a+bi$ explicitly and compute both sides of the equation step by step.
    \item \textbf{Penultimate Step}: Once we have values for $a$ and $b$, the final step is computing $\sqrt{a^2+b^2}$.
    \item \textbf{Make It Easier}: Consider what happens with specific simple cases (e.g., purely real or purely imaginary $z$) to build intuition, though the full algebraic approach is most direct.
    \item \textbf{Conjecture via Small Cases}: The equation $z \otimes z = z^2 + 40$ suggests that the operation $\otimes$ differs from standard squaring by a constant 40.
    \item \textbf{Visualization}: Complex numbers on the plane; we're looking for the distance from origin to point $(a,b)$.
\end{itemize}

\subsection*{C. Argument Construction}
\begin{itemize}[leftmargin=*]
    \item \textbf{Proof Method}: Direct computation - set up equations and solve algebraically
    \item \textbf{Key Lemmas}: 
    \begin{enumerate}
        \item For $z = a+bi$: $z \otimes z = a^2 + b^2i$
        \item For $z = a+bi$: $z^2 = a^2 - b^2 + 2abi$
        \item Equating real parts: $a^2 = a^2 - b^2 + 40$
        \item Equating imaginary parts: $b^2 = 2ab$
    \end{enumerate}
    \item \textbf{Logical Structure}: Write $z$ in standard form $\to$ compute both sides $\to$ equate components $\to$ solve system $\to$ compute magnitude
\end{itemize}
\end{strategicbox}

\begin{tacticalbox}
\subsection*{A. Core Mathematical Tactics}
\begin{itemize}[leftmargin=*]
    \item \textbf{Primary Tactics}: 
    \begin{itemize}
        \item \emph{Algebraic Manipulation}: Expand and simplify complex expressions
        \item \emph{Equating Real and Imaginary Parts}: Standard technique for complex number equations
        \item \emph{System Solving}: Solve coupled equations for $a$ and $b$
    \end{itemize}
    \item \textbf{Domain-Specific Tactics (Complex Numbers)}: 
    \begin{itemize}
        \item Write complex numbers in rectangular form $a+bi$
        \item Separate real and imaginary components
        \item Apply magnitude formula $|z| = \sqrt{a^2+b^2}$
        \item Understand custom operations by applying definition carefully
    \end{itemize}
    \item \textbf{Crossover Tactics}: This problem bridges algebra and complex analysis, requiring careful symbolic manipulation.
\end{itemize}
\end{tacticalbox}

\begin{techniquesbox}
\subsection*{A. High-Probability Techniques}
\begin{itemize}[leftmargin=*]
    \item \textbf{Primary Technique}: \textbf{Complex Number Conversion (H18)} - Converting to rectangular form $a+bi$ and separating real/imaginary parts
    \item \textbf{Recognition Cues}: 
    \begin{itemize}
        \item Problem involves complex numbers with custom operation
        \item Need to find magnitude $|z|$
        \item Equation involves complex number equality (equate both components)
    \end{itemize}
    \item \textbf{Application}: 
    \begin{enumerate}
        \item Let $z = a+bi$
        \item Compute $z \otimes z = a \cdot a + b \cdot b \cdot i = a^2 + b^2i$
        \item Compute $z^2 = (a+bi)^2 = a^2 + 2abi + (bi)^2 = a^2 - b^2 + 2abi$
        \item Set up equation: $a^2 + b^2i = (a^2 - b^2 + 2abi) + 40$
        \item Simplify: $a^2 + b^2i = (a^2 - b^2 + 40) + 2abi$
        \item Equate real parts: $a^2 = a^2 - b^2 + 40 \Rightarrow b^2 = 40$
        \item Equate imaginary parts: $b^2 = 2ab \Rightarrow b = 2a$ (assuming $b \neq 0$)
    \end{enumerate}
\end{itemize}

\subsection*{B. Alternative Techniques}
\begin{itemize}[leftmargin=*]
    \item \textbf{Backup Techniques}: 
    \begin{itemize}
        \item \textbf{Answer Choice Testing}: Plug in answer choices - if $|z| = k$, then $a^2 + b^2 = k^2$. Use constraints $b^2 = 40$ and $b = 2a$ to check which value of $k$ works.
        \item \textbf{Working Backwards (H6)}: From answer choices, determine what $a^2+b^2$ should equal, then verify consistency with the given equation.
    \end{itemize}
    \item \textbf{Technique Integration}: Primary algebraic approach is most direct. Answer choice verification serves as excellent verification method.
\end{itemize}
\end{techniquesbox}

\begin{verificationbox}
\subsection*{A. Verification Level Selection}
\begin{itemize}[leftmargin=*]
    \item \textbf{Chosen Level}: Level 4 (Standard Verification, 30-60 seconds)
    \item \textbf{Justification}: Mid-band problem (P12) with algebraic manipulation that's prone to sign errors and computational mistakes. Standard verification is appropriate to catch these errors.
    \item \textbf{Time Budget}: 45 seconds for verification
\end{itemize}

\subsection*{B. Verification Checklist}
\begin{itemize}[leftmargin=*]
    \item \textbf{Domain Restrictions}: Check that $b \neq 0$ (which is satisfied since $b^2 = 40$)
    \item \textbf{Computational Accuracy}: 
    \begin{itemize}
        \item Verify $b^2 = 40 \Rightarrow b = \pm 2\sqrt{10}$
        \item Verify $b = 2a \Rightarrow a = \pm \sqrt{10}$
        \item Verify $a^2 + b^2 = 10 + 40 = 50$
        \item Verify $|z| = \sqrt{50} = 5\sqrt{2}$
    \end{itemize}
    \item \textbf{Alternative Method}: Substitute back into original equation $z \otimes z = z^2 + 40$ to verify consistency
    \item \textbf{Edge Cases}: Check that both equations (real and imaginary parts) are satisfied simultaneously
    \item \textbf{Answer Choice Check}: Verify answer matches one of the given choices (E)
\end{itemize}

\subsection*{C. Detailed Verification}
\begin{itemize}[leftmargin=*]
    \item \textbf{Method 1: Direct Substitution}
    
    With $a = \sqrt{10}$, $b = 2\sqrt{10}$ (taking positive values):
    \begin{align*}
        z &= \sqrt{10} + 2\sqrt{10}i\\
        z \otimes z &= (\sqrt{10})^2 + (2\sqrt{10})^2 i = 10 + 40i\\
        z^2 &= (\sqrt{10})^2 - (2\sqrt{10})^2 + 2(\sqrt{10})(2\sqrt{10})i\\
        &= 10 - 40 + 40i = -30 + 40i\\
        z^2 + 40 &= -30 + 40 + 40i = 10 + 40i \checkmark
    \end{align*}
    
    \item \textbf{Method 2: Answer Choice Verification}
    
    Testing $|z| = 5\sqrt{2}$:
    \begin{align*}
        |z|^2 &= (5\sqrt{2})^2 = 25 \cdot 2 = 50\\
        a^2 + b^2 &= 50
    \end{align*}
    With $b^2 = 40$, we get $a^2 = 10$, so $a = \pm\sqrt{10}$.
    
    Check $b = 2a$: $b^2 = 4a^2 \Rightarrow 40 = 4(10) = 40$ \checkmark
\end{itemize}
\end{verificationbox}

\begin{solutionbox}
\subsection*{A. Step-by-Step Solution}
\begin{enumerate}[leftmargin=*]
    \item \textbf{Initial Setup (30 seconds)}
    
    Let $z = a + bi$ where $a, b \in \mathbb{R}$.
    
    Our goal is to find $|z| = \sqrt{a^2 + b^2}$.
    
    \item \textbf{Compute $z \otimes z$ (30 seconds)}
    
    Using the definition $u \otimes v = ac + bdi$ with $u = v = z = a+bi$:
    \begin{align*}
        z \otimes z &= (a)(a) + (b)(b)i\\
        &= a^2 + b^2i
    \end{align*}
    
    \item \textbf{Compute $z^2$ (30 seconds)}
    
    Using standard complex multiplication:
    \begin{align*}
        z^2 &= (a+bi)^2\\
        &= a^2 + 2abi + (bi)^2\\
        &= a^2 + 2abi + b^2i^2\\
        &= a^2 + 2abi - b^2\\
        &= (a^2 - b^2) + 2abi
    \end{align*}
    
    \item \textbf{Set up the equation (15 seconds)}
    
    Given: $z \otimes z = z^2 + 40$
    
    Substituting:
    $$a^2 + b^2i = (a^2 - b^2 + 40) + 2abi$$
    
    \item \textbf{Equate real parts (45 seconds)}
    
    Real part of LHS: $a^2$
    
    Real part of RHS: $a^2 - b^2 + 40$
    
    Equation: $a^2 = a^2 - b^2 + 40$
    
    Simplifying: $0 = -b^2 + 40$
    
    Therefore: $b^2 = 40$, so $b = \pm 2\sqrt{10}$
    
    \item \textbf{Equate imaginary parts (45 seconds)}
    
    Imaginary coefficient of LHS: $b^2$
    
    Imaginary coefficient of RHS: $2ab$
    
    Equation: $b^2 = 2ab$
    
    Since $b \neq 0$ (we found $b^2 = 40$), we can divide by $b$:
    $$b = 2a$$
    
    Therefore: $a = \frac{b}{2} = \frac{\pm 2\sqrt{10}}{2} = \pm\sqrt{10}$
    
    \item \textbf{Compute magnitude (30 seconds)}
    
    \begin{align*}
        |z| &= \sqrt{a^2 + b^2}\\
        &= \sqrt{(\pm\sqrt{10})^2 + (\pm 2\sqrt{10})^2}\\
        &= \sqrt{10 + 4 \cdot 10}\\
        &= \sqrt{10 + 40}\\
        &= \sqrt{50}\\
        &= \sqrt{25 \cdot 2}\\
        &= 5\sqrt{2}
    \end{align*}
\end{enumerate}

\subsection*{B. Alternative Approaches}

\textbf{Method 2: Answer Choice Testing}

Since we have multiple choice answers, we can work backwards:

\begin{itemize}
    \item If $|z| = 5\sqrt{2}$, then $a^2 + b^2 = 50$
    \item From equating real parts: $b^2 = 40$
    \item Therefore: $a^2 = 50 - 40 = 10$, so $a = \pm\sqrt{10}$
    \item Check imaginary equation: $b = 2a \Rightarrow b^2 = 4a^2$
    \item Verify: $40 = 4(10) = 40$ \checkmark
\end{itemize}

This confirms $|z| = 5\sqrt{2}$.

\textbf{Method 3: Pattern Recognition}

Notice that the answer $5\sqrt{2} = \sqrt{50}$ and we have the constant 40 in the problem. This suggests looking for $a^2 + b^2 = 50$. Working backwards from this insight, we can quickly verify it satisfies all conditions.
\end{solutionbox}

\begin{competitionbox}
\subsection*{A. Time Management}
\begin{itemize}[leftmargin=*]
    \item \textbf{Estimated Time}: 3.5-4 minutes total
    \begin{itemize}
        \item Setup and computation: 2.5-3 minutes
        \item Verification: 45-60 seconds
    \end{itemize}
    \item \textbf{Time Checkpoints}: 
    \begin{itemize}
        \item At 1.5 minutes: Should have computed both $z \otimes z$ and $z^2$
        \item At 2.5 minutes: Should have solved for $b^2$ and be working on finding $a$
        \item If stuck at 3 minutes: Consider answer choice testing method
    \end{itemize}
    \item \textbf{Priority Level}: High - This is a mid-band problem that should be completed for students aiming for scores of 90+
\end{itemize}

\subsection*{B. Risk Assessment}
\begin{itemize}[leftmargin=*]
    \item \textbf{Confidence Level}: 85-90\% after standard verification
    \item \textbf{Abandonment Criteria}: 
    \begin{itemize}
        \item If computation becomes very messy after 4 minutes, mark for review and move on
        \item However, this problem should be straightforward with careful algebra
    \end{itemize}
    \item \textbf{Flag for Review}: Only if arithmetic errors are suspected. Otherwise, move on confidently after verification.
\end{itemize}

\subsection*{C. Competition Context}
\begin{itemize}[leftmargin=*]
    \item \textbf{Problem 12/25 Strategy}: This falls in the "must get" range for competitive students. Allocate full 3-4 minutes without rushing.
    \item \textbf{Point Value}: 6 points (standard AMC12), with 1.5 points for leaving blank
    \item \textbf{Difficulty Assessment}: Standard mid-band problem. Tests:
    \begin{itemize}
        \item Understanding of complex number operations
        \item Careful reading (the operation $\otimes$ is non-standard)
        \item Algebraic manipulation skills
        \item System solving ability
    \end{itemize}
\end{itemize}
\end{competitionbox}

\begin{keyinsightbox}
\textbf{Key Insight}: The custom operation $u \otimes v = ac + bdi$ is NOT standard complex multiplication. It extracts the real part from the first factor's product ($ac$) and the imaginary part from the second factor's product ($bdi$). This creates a different structure when computing $z \otimes z = a^2 + b^2i$, which leads to the equation $b^2 = 40$ from the real parts and $b = 2a$ from the imaginary parts.

The problem elegantly combines non-standard definitions with standard complex number theory, rewarding careful reading and systematic algebraic manipulation.
\end{keyinsightbox}

\begin{warningbox}
\textbf{Warning}: Common Pitfalls to Avoid

\begin{itemize}
    \item \textbf{Misapplying the operation}: Do NOT compute $z \otimes z$ as if it were standard complex multiplication. The operation is specifically defined as $ac + bdi$, which differs from $(a+bi)(c+di) = (ac-bd) + (ad+bc)i$.
    
    \item \textbf{Sign errors}: When computing $z^2 = (a+bi)^2$, remember that $(bi)^2 = b^2i^2 = -b^2$, not $+b^2$.
    
    \item \textbf{Forgetting absolute value}: When solving $b = 2a$ with $b^2 = 40$, remember that both $a$ and $b$ can be positive or negative. However, the magnitude $|z| = \sqrt{a^2+b^2}$ is always positive and independent of the signs.
    
    \item \textbf{Dividing by zero}: When simplifying $b^2 = 2ab$ to $b = 2a$, we must verify that $b \neq 0$. This is confirmed since $b^2 = 40 > 0$.
    
    \item \textbf{Computational errors}: $\sqrt{50} = \sqrt{25 \cdot 2} = 5\sqrt{2}$, not $5\sqrt{10}$ or other variations.
\end{itemize}
\end{warningbox}

\begin{tipbox}
\textbf{Pro Tip}: Competition Strategy

\begin{itemize}
    \item \textbf{Read carefully}: Non-standard operations are common in AMC problems. Always verify what operation is actually defined before computing.
    
    \item \textbf{Use answer choices}: The presence of $\sqrt{10}$ and $5\sqrt{2} = \sqrt{50}$ in the choices hints that the answer involves these values. The constant 40 in the problem strongly suggests $b^2 = 40$.
    
    \item \textbf{Separate components}: For any complex number equation, equating real and imaginary parts separately is a fundamental technique. This transforms one complex equation into two real equations.
    
    \item \textbf{Verify systematically}: After finding $a$ and $b$, substitute back into the original equation to ensure both sides are equal. This catches sign errors and computational mistakes.
    
    \item \textbf{Time management}: This is a standard mid-band problem. Don't overthink it - follow the systematic approach and you should complete it in 3-4 minutes.
\end{itemize}
\end{tipbox}

\section*{Final Answer}

The magnitude of $z$ is:
$$\boxed{\textbf{(E) } 5\sqrt{2}}$$

\section*{Confidence Assessment}
\begin{itemize}[leftmargin=*]
    \item \textbf{Confidence Level}: 90\% - The solution is straightforward with careful algebra, and verification by substitution confirms correctness.
    \item \textbf{Verification Applied}: Level 4 (Standard Verification) - Checked by direct substitution and answer choice verification
    \item \textbf{Flag for Review}: No - Solution is complete and verified through multiple methods
    \item \textbf{Time Spent}: Estimated 3.5-4 minutes (within target for P12)
\end{itemize}

\section*{Summary of Solution Path}

\begin{enumerate}
    \item Let $z = a + bi$ and apply the custom operation: $z \otimes z = a^2 + b^2i$
    \item Compute standard complex square: $z^2 = (a^2-b^2) + 2abi$
    \item Set up equation: $a^2 + b^2i = (a^2-b^2+40) + 2abi$
    \item Equate real parts: $a^2 = a^2 - b^2 + 40 \Rightarrow b^2 = 40$
    \item Equate imaginary parts: $b^2 = 2ab \Rightarrow b = 2a \Rightarrow a^2 = 10$
    \item Compute magnitude: $|z| = \sqrt{a^2+b^2} = \sqrt{10+40} = \sqrt{50} = 5\sqrt{2}$
\end{enumerate}

\end{document}

